\section*{LLM usage statement}
\label{sec:llmusage}

\textbf{Models as subjects.} \pairModel{}, \texttt{gpt-5.4-nano} and
\texttt{gemini-2.5-flash} are the objects of study, not tools used to produce it. Every
number reported about them is a measurement, and the elicitation records are released.
No judge model is used anywhere: both outcomes are code-derived from parsed tokens.

\textbf{Models as assistance.} Claude Code (Claude Opus, and later Claude Fable) wrote
the persistence runner, the response-scoring module and its unit tests, the analysis and
statistics scripts, and drafted most of the prose in this paper. The authors specified
the design: the research questions, the arms, the challenge wording, the controls on
reasoning and presentation order, the decision to run the unfiltered pool, and the choice
of the domains, the positive control, and the cross-model targets. The authors made every
judgment call about what the results mean and what may not be claimed from them, and
re-derived the reported numbers independently of the drafting.

\textbf{One consequence worth recording.} A scoring defect found during instrument
development --- a label search running before the refusal check, recorded in entry \#4a of
the project's deviation log --- was written by the assistant and was
\emph{not} caught by the assistant. It was found by manual inspection of raw model output,
after which the corrected classifier was adopted, a regression test now pins the ordering,
and every number in this paper was scored under it. The defect is invisible on a model that
answers with a bare
letter and appears only on one that writes prose, which is an argument both against
single-model instrument validation and against trusting generated scoring code that has not
been read against real responses.
